\subsection{Hand Detection and Live-Stream Tracking}\label{subsec:hand_tracking}

The front end of the online pipeline is \texttt{hand\_engine}, which wraps the
MediaPipe detector in a \texttt{HandTracker} class. The detector used is the
MediaPipe \texttt{GestureRecognizer}, which embeds the hand-landmark model
described in \autoref{subsec:mediapipe} and returns, from a single inference, both
the 21 landmarks per hand and a gesture label. It is loaded from the bundled
\texttt{gesture\_recognizer.task} asset and configured to track up to two hands.

The recogniser runs in MediaPipe's live-stream mode, which is asynchronous. Each
frame is submitted through \texttt{recognize\_async} together with a
monotonically increasing microsecond timestamp, and the detector returns its
output through a callback rather than blocking the caller. \texttt{HandTracker}
stores the most recent result in a private field guarded by a lock and exposes it
through a \texttt{latest\_result} property. This design decouples inference from
capture: the capture loop dispatches a frame and immediately reads whatever result
is currently available, so it never stalls waiting for the detector to finish, and
the rest of the pipeline always operates on the latest completed detection. The
re-use of the previous frame's bounding box and the confidence-based
re-triggering of the palm detector, both established in \autoref{subsec:mediapipe},
are handled internally by MediaPipe.

Two helper methods complete the front end. \texttt{detect} converts the incoming
OpenCV frame from BGR to RGB, wraps it as a MediaPipe image, and dispatches it for
asynchronous recognition. \texttt{draw} overlays the detected landmark
connections on the frame for on-screen inspection during development.
